
\documentclass{article}
\usepackage{colortbl}
\usepackage{makecell}
\usepackage{multirow}
\usepackage{supertabular}

\begin{document}

\newcounter{utterance}

\centering \large Interaction Transcript for game `function\_detective', experiment `two\_numbers\_passive\_labeled\_pairs', episode 261 with Qwen3.6{-}35B{-}A3B{-}FP8{-}without{-}reasoning.
\vspace{24pt}

{ \footnotesize  \setcounter{utterance}{1}
\setlength{\tabcolsep}{0pt}
\begin{supertabular}{c@{$\;$}|p{.15\linewidth}@{}p{.15\linewidth}p{.15\linewidth}p{.15\linewidth}p{.15\linewidth}p{.15\linewidth}}
    \# & \multicolumn{2}{c}{Player} && \multicolumn{2}{c}{Game Master} \\
    \hline

    \theutterance \stepcounter{utterance}  
    & & & \multicolumn{4}{p{0.6\linewidth}}{
        \cellcolor[rgb]{0.9,0.9,0.9}{
            \makecell[{{p{\linewidth}}}]{
                \texttt{\tiny{[P1$\langle$GM]}}
                \texttt{I have a hidden Python function for you to figure out.} \\
\\ 
\texttt{Your job is to infer what it does and then submit a correct Python implementation.} \\
\\ 
\texttt{Function signature:} \\
\texttt{(a: int, b: int) {-}> int} \\
\\ 
\texttt{You have 15 turns total.} \\
\\ 
\texttt{I reveal candidate input{-}output pairs one at a time and tell you whether each one is correct.} \\
\\ 
\texttt{You must make your best SOLVE attempt no later than the second{-}to{-}last turn.} \\
\texttt{Do not wait until the very end.} \\
\\ 
\texttt{You can use these actions:} \\
\texttt{{-} JSON format described at the end of the prompt} \\
\texttt{{-} SOLVE: ```python ... ```} \\
\\ 
\texttt{Do not add import statements in your guess.} \\
\texttt{Anything you might need is already available, so just write the function body in hypothesis and return the answer.} \\
\\ 
\texttt{Preloaded examples:} \\
\texttt{1. (((0, {-}2), {-}2), True)} \\
\\ 
\texttt{Scoring:} \\
\texttt{{-} correctness is required for any non{-}zero score} \\
\texttt{{-} in interactive modes, solving with fewer information requests is better} \\
\\ 
\texttt{Input formatting reminders:} \\
\texttt{{-} strings must be quoted} \\
\texttt{{-} booleans must be True/False} \\
\texttt{{-} use commas between multiple values} \\
\texttt{{-} do not write any extra action tags in the same message} \\
\\ 
\texttt{Format example:} \\
\texttt{You write:} \\
\texttt{\{'current\_mode': <probing/confirming/don't know; choose whichever best describes your current state in the task>, 'input\_rationale': <brief description of why you chose this hypothesis>, 'hypothesis': <either the exact string "unknown" if you have no guess yet OR your current best guess as runnable Python matching the given signature, with a one{-}line docstring stating what it does. Never write placeholder code, if you cannot commit to a guess, use "unknown". Match the given signature.} \\
\texttt{e.g.} \\
\texttt{def f(x):} \\
\texttt{	"""Returns the absolute value of x."""} \\
\texttt{	return x if x >= 0 else {-}x} \\
\\ 
\texttt{Use "unknown" if you have no guess.>, 'confidence\_score': <score between 1{-}100. 100 means highest confidence in your hypothesis>\}} \\
\\ 
\texttt{Output only one valid json per turn.} \\
            }
        }
    }
     \\ \\

    \theutterance \stepcounter{utterance}  
    & \multicolumn{4}{p{0.6\linewidth}}{
        \cellcolor[rgb]{0.561,0.737,0.561}{
            \makecell[{{p{\linewidth}}}]{
                \texttt{\tiny{[P1$\rangle$GM]}}
                \texttt{\{'current\_mode': 'probing', 'input\_rationale': 'I have one preloaded example: f(0, {-}2) = True. However, the function signature specifies it returns an int, not a bool. This suggests the example might be formatted in a way where True is treated as 1, or the example is testing a boolean condition that maps to an integer. Let's probe with simple inputs to understand the behavior. I will try f(0, 0) and f(1, 0) to see if it's related to GCD, absolute difference, or some other common integer function.', 'hypothesis': 'unknown', 'confidence\_score': 10\}} \\
            }
        }
    }
    & & \\ \\

    \theutterance \stepcounter{utterance}  
    & & & \multicolumn{4}{p{0.6\linewidth}}{
        \cellcolor[rgb]{0.9,0.9,0.9}{
            \makecell[{{p{\linewidth}}}]{
                \texttt{\tiny{[P1$\langle$GM]}}
                \texttt{Game Violation: ParseError: Response is not valid JSON.} \\
\texttt{Output ONLY:} \\
\texttt{One JSON with confidence\_score, current\_mode, input\_rationale, and hypothesis keys ...} \\
\texttt{OR} \\
\texttt{SOLVE: ```python ... ```} \\
            }
        }
    }
     \\ \\

    \theutterance \stepcounter{utterance}  
    & \multicolumn{4}{p{0.6\linewidth}}{
        \cellcolor[rgb]{0.561,0.737,0.561}{
            \makecell[{{p{\linewidth}}}]{
                \texttt{\tiny{[P1$\rangle$GM]}}
                \texttt{\{} \\
\texttt{  "current\_mode": "probing",} \\
\texttt{  "input\_rationale": "The preloaded example (((0, {-}2), {-}2), True) is ambiguous. It might mean f(0, {-}2) returns True (which is 1 in Python int context) or it's a nested structure. Given the signature (a: int, b: int) {-}> int, I suspect the 'True' is actually the integer 1. Let's test simple cases to distinguish between common functions like GCD, absolute difference, or simple arithmetic. I'll start with f(0, 0) and f(1, 1).",} \\
\texttt{  "hypothesis": "unknown",} \\
\texttt{  "confidence\_score": 10} \\
\texttt{\}} \\
            }
        }
    }
    & & \\ \\

    \theutterance \stepcounter{utterance}  
    & & & \multicolumn{4}{p{0.6\linewidth}}{
        \cellcolor[rgb]{0.9,0.9,0.9}{
            \makecell[{{p{\linewidth}}}]{
                \texttt{\tiny{[P1$\langle$GM]}}
                \texttt{OUTPUT: ((({-}1, 99), 1), False)} \\
            }
        }
    }
     \\ \\

    \theutterance \stepcounter{utterance}  
    & \multicolumn{4}{p{0.6\linewidth}}{
        \cellcolor[rgb]{0.561,0.737,0.561}{
            \makecell[{{p{\linewidth}}}]{
                \texttt{\tiny{[P1$\rangle$GM]}}
                \texttt{\{} \\
\texttt{  "current\_mode": "probing",} \\
\texttt{  "input\_rationale": "I have two data points now:\textbackslash{}n1. f(0, {-}2) = True (which is 1 in integer context)\textbackslash{}n2. f({-}1, 99) = False (which is 0 in integer context)\textbackslash{}n\textbackslash{}nLet's analyze:\textbackslash{}n{-} f(0, {-}2) = 1\textbackslash{}n{-} f({-}1, 99) = 0\textbackslash{}n\textbackslash{}nThis looks like it could be related to whether the two numbers are coprime (GCD is 1). Let's check:\textbackslash{}n{-} GCD(0, {-}2) = 2. Not 1. So not simple coprimality.\textbackslash{}n{-} Maybe it's checking if a + b is even? 0 + ({-}2) = {-}2 (even). {-}1 + 99 = 98 (even). Both even, but results differ. So not parity of sum.\textbackslash{}n{-} Maybe it's checking if a and b have the same sign? 0 is neither positive nor negative. {-}1 is negative, 99 is positive. Different signs {-}> 0. 0 and {-}2... 0 is tricky.\textbackslash{}n{-} Let's try f(1, 1). If it's coprimality, GCD(1,1)=1 {-}> True/1. If it's same sign, 1 and 1 are same} \\
            }
        }
    }
    & & \\ \\

    \theutterance \stepcounter{utterance}  
    & & & \multicolumn{4}{p{0.6\linewidth}}{
        \cellcolor[rgb]{0.9,0.9,0.9}{
            \makecell[{{p{\linewidth}}}]{
                \texttt{\tiny{[P1$\langle$GM]}}
                \texttt{Game Violation: ParseError: Response is not valid JSON.} \\
\texttt{Output ONLY:} \\
\texttt{One JSON with confidence\_score, current\_mode, input\_rationale, and hypothesis keys ...} \\
\texttt{OR} \\
\texttt{SOLVE: ```python ... ```} \\
            }
        }
    }
     \\ \\

    \theutterance \stepcounter{utterance}  
    & \multicolumn{4}{p{0.6\linewidth}}{
        \cellcolor[rgb]{0.561,0.737,0.561}{
            \makecell[{{p{\linewidth}}}]{
                \texttt{\tiny{[P1$\rangle$GM]}}
                \texttt{\{} \\
\texttt{  "current\_mode": "probing",} \\
\texttt{  "input\_rationale": "I have two data points: f(0, {-}2) = 1 (True) and f({-}1, 99) = 0 (False). The function returns an int (0 or 1). Let's test more cases to narrow down the logic. Common candidates: GCD(a,b)==1, a==b, a+b==0, etc. f(0,{-}2)=1 suggests GCD(0,{-}2)=2 != 1, so not GCD==1. f({-}1,99)=0. Let's try f(1,1) and f(2,2) to check equality or GCD.",} \\
\texttt{  "hypothesis": "unknown",} \\
\texttt{  "confidence\_score": 20} \\
\texttt{\}} \\
            }
        }
    }
    & & \\ \\

    \theutterance \stepcounter{utterance}  
    & & & \multicolumn{4}{p{0.6\linewidth}}{
        \cellcolor[rgb]{0.9,0.9,0.9}{
            \makecell[{{p{\linewidth}}}]{
                \texttt{\tiny{[P1$\langle$GM]}}
                \texttt{OUTPUT: (((10, {-}1000), {-}1000), True)} \\
            }
        }
    }
     \\ \\

    \theutterance \stepcounter{utterance}  
    & \multicolumn{4}{p{0.6\linewidth}}{
        \cellcolor[rgb]{0.561,0.737,0.561}{
            \makecell[{{p{\linewidth}}}]{
                \texttt{\tiny{[P1$\rangle$GM]}}
                \texttt{\{} \\
\texttt{  "current\_mode": "probing",} \\
\texttt{  "input\_rationale": "I have three data points:\textbackslash{}n1. f(0, {-}2) = 1 (True)\textbackslash{}n2. f({-}1, 99) = 0 (False)\textbackslash{}n3. f(10, {-}1000) = 1 (True)\textbackslash{}n\textbackslash{}nLet's look at the relationship between inputs and outputs.\textbackslash{}n{-} f(0, {-}2) = 1. Note that 0 is a multiple of {-}2 (0 = 0 * {-}2). Or {-}2 is a multiple of 0? No. But 0 is divisible by {-}2.\textbackslash{}n{-} f({-}1, 99) = 0. {-}1 is not a multiple of 99, and 99 is not a multiple of {-}1? Wait, 99 = {-}99 * {-}1. So 99 IS a multiple of {-}1. But the result is 0. So it's not simply 'is a multiple of'.\textbackslash{}n{-} f(10, {-}1000) = 1. {-}1000 = {-}100 * 10. So {-}1000 is a multiple of 10. Result is 1.\textbackslash{}n\textbackslash{}nLet's re{-}evaluate f({-}1, 99). Is {-}1 a multiple of 99?} \\
            }
        }
    }
    & & \\ \\

    \theutterance \stepcounter{utterance}  
    & & & \multicolumn{4}{p{0.6\linewidth}}{
        \cellcolor[rgb]{0.9,0.9,0.9}{
            \makecell[{{p{\linewidth}}}]{
                \texttt{\tiny{[P1$\langle$GM]}}
                \texttt{Game Violation: ParseError: Response is not valid JSON.} \\
\texttt{Output ONLY:} \\
\texttt{One JSON with confidence\_score, current\_mode, input\_rationale, and hypothesis keys ...} \\
\texttt{OR} \\
\texttt{SOLVE: ```python ... ```} \\
            }
        }
    }
     \\ \\

    \theutterance \stepcounter{utterance}  
    & \multicolumn{4}{p{0.6\linewidth}}{
        \cellcolor[rgb]{0.561,0.737,0.561}{
            \makecell[{{p{\linewidth}}}]{
                \texttt{\tiny{[P1$\rangle$GM]}}
                \texttt{\{} \\
\texttt{  "current\_mode": "probing",} \\
\texttt{  "input\_rationale": "I have three data points: f(0, {-}2) = 1, f({-}1, 99) = 0, f(10, {-}1000) = 1. The pattern seems to be related to divisibility. Specifically, it looks like it checks if the second argument divides the first argument (a \% b == 0). Let's verify: 0 \% {-}2 == 0 (True). 10 \% {-}1000 == 10 != 0 (False? Wait, 10 is not divisible by {-}1000). But the result is True. So it's not a \% b == 0. Maybe it's b \% a == 0? {-}2 \% 0 is undefined. So that's not it. Let's try checking if a divides b or b divides a. 0 is divisible by {-}2. 10 is NOT divisible by {-}1000. {-}1000 IS divisible by 10. So f(10, {-}1000) checks if {-}1000 is divisible by 10. That is True. f(0, {-}2) checks if {-}2 is divisible by 0? No. It checks if 0 is divisible by {-}2. That is True. f({-}1,} \\
            }
        }
    }
    & & \\ \\

    \theutterance \stepcounter{utterance}  
    & & & \multicolumn{4}{p{0.6\linewidth}}{
        \cellcolor[rgb]{0.9,0.9,0.9}{
            \makecell[{{p{\linewidth}}}]{
                \texttt{\tiny{[P1$\langle$GM]}}
                \texttt{Game Violation: ParseError: Response is not valid JSON.} \\
\texttt{Output ONLY:} \\
\texttt{One JSON with confidence\_score, current\_mode, input\_rationale, and hypothesis keys ...} \\
\texttt{OR} \\
\texttt{SOLVE: ```python ... ```} \\
            }
        }
    }
     \\ \\

    \theutterance \stepcounter{utterance}  
    & \multicolumn{4}{p{0.6\linewidth}}{
        \cellcolor[rgb]{0.561,0.737,0.561}{
            \makecell[{{p{\linewidth}}}]{
                \texttt{\tiny{[P1$\rangle$GM]}}
                \texttt{\{} \\
\texttt{  "current\_mode": "probing",} \\
\texttt{  "input\_rationale": "I have three data points: f(0, {-}2) = 1, f({-}1, 99) = 0, f(10, {-}1000) = 1. The pattern seems to be related to divisibility. Specifically, it looks like it checks if the second argument divides the first argument (a \% b == 0). Let's verify: 0 \% {-}2 == 0 (True). 10 \% {-}1000 == 10 != 0 (False? Wait, 10 is not divisible by {-}1000). But the result is True. So it's not a \% b == 0. Maybe it's b \% a == 0? {-}2 \% 0 is undefined. So that's not it. Let's try checking if a divides b or b divides a. 0 is divisible by {-}2. 10 is NOT divisible by {-}1000. {-}1000 IS divisible by 10. So f(10, {-}1000) checks if {-}1000 is divisible by 10. That is True. f(0, {-}2) checks if {-}2 is divisible by 0? No. It checks if 0 is divisible by {-}2. That is True. f({-}1,} \\
            }
        }
    }
    & & \\ \\

    \theutterance \stepcounter{utterance}  
    & & & \multicolumn{4}{p{0.6\linewidth}}{
        \cellcolor[rgb]{0.9,0.9,0.9}{
            \makecell[{{p{\linewidth}}}]{
                \texttt{\tiny{[P1$\langle$GM]}}
                \texttt{Game Violation: ParseError: Response is not valid JSON.} \\
\texttt{Output ONLY:} \\
\texttt{One JSON with confidence\_score, current\_mode, input\_rationale, and hypothesis keys ...} \\
\texttt{OR} \\
\texttt{SOLVE: ```python ... ```} \\
            }
        }
    }
     \\ \\

    \theutterance \stepcounter{utterance}  
    & \multicolumn{4}{p{0.6\linewidth}}{
        \cellcolor[rgb]{0.561,0.737,0.561}{
            \makecell[{{p{\linewidth}}}]{
                \texttt{\tiny{[P1$\rangle$GM]}}
                \texttt{\{} \\
\texttt{  "current\_mode": "probing",} \\
\texttt{  "input\_rationale": "I have three data points: f(0, {-}2) = 1, f({-}1, 99) = 0, f(10, {-}1000) = 1. The pattern seems to be related to divisibility. Specifically, it looks like it checks if the second argument divides the first argument (a \% b == 0). Let's verify: 0 \% {-}2 == 0 (True). 10 \% {-}1000 == 10 != 0 (False? Wait, 10 is not divisible by {-}1000). But the result is True. So it's not a \% b == 0. Maybe it's b \% a == 0? {-}2 \% 0 is undefined. So that's not it. Let's try checking if a divides b or b divides a. 0 is divisible by {-}2. 10 is NOT divisible by {-}1000. {-}1000 IS divisible by 10. So f(10, {-}1000) checks if {-}1000 is divisible by 10. That is True. f(0, {-}2) checks if {-}2 is divisible by 0? No. It checks if 0 is divisible by {-}2. That is True. f({-}1,} \\
            }
        }
    }
    & & \\ \\

    \theutterance \stepcounter{utterance}  
    & & & \multicolumn{4}{p{0.6\linewidth}}{
        \cellcolor[rgb]{0.9,0.9,0.9}{
            \makecell[{{p{\linewidth}}}]{
                \texttt{\tiny{[P1$\langle$GM]}}
                \texttt{Game Violation: ParseError: Response is not valid JSON.} \\
\texttt{Output ONLY:} \\
\texttt{One JSON with confidence\_score, current\_mode, input\_rationale, and hypothesis keys ...} \\
\texttt{OR} \\
\texttt{SOLVE: ```python ... ```} \\
            }
        }
    }
     \\ \\

    \theutterance \stepcounter{utterance}  
    & \multicolumn{4}{p{0.6\linewidth}}{
        \cellcolor[rgb]{0.561,0.737,0.561}{
            \makecell[{{p{\linewidth}}}]{
                \texttt{\tiny{[P1$\rangle$GM]}}
                \texttt{\{} \\
\texttt{  "current\_mode": "probing",} \\
\texttt{  "input\_rationale": "I have three data points: f(0, {-}2) = 1, f({-}1, 99) = 0, f(10, {-}1000) = 1. The pattern seems to be related to divisibility. Specifically, it looks like it checks if the second argument divides the first argument (a \% b == 0). Let's verify: 0 \% {-}2 == 0 (True). 10 \% {-}1000 == 10 != 0 (False? Wait, 10 is not divisible by {-}1000). But the result is True. So it's not a \% b == 0. Maybe it's b \% a == 0? {-}2 \% 0 is undefined. So that's not it. Let's try checking if a divides b or b divides a. 0 is divisible by {-}2. 10 is NOT divisible by {-}1000. {-}1000 IS divisible by 10. So f(10, {-}1000) checks if {-}1000 is divisible by 10. That is True. f(0, {-}2) checks if {-}2 is divisible by 0? No. It checks if 0 is divisible by {-}2. That is True. f({-}1,} \\
            }
        }
    }
    & & \\ \\

    \theutterance \stepcounter{utterance}  
    & & & \multicolumn{4}{p{0.6\linewidth}}{
        \cellcolor[rgb]{0.9,0.9,0.9}{
            \makecell[{{p{\linewidth}}}]{
                \texttt{\tiny{[P1$\langle$GM]}}
                \texttt{Game Violation: ParseError: Response is not valid JSON.} \\
\texttt{Output ONLY:} \\
\texttt{One JSON with confidence\_score, current\_mode, input\_rationale, and hypothesis keys ...} \\
\texttt{OR} \\
\texttt{SOLVE: ```python ... ```} \\
            }
        }
    }
     \\ \\

    \theutterance \stepcounter{utterance}  
    & \multicolumn{4}{p{0.6\linewidth}}{
        \cellcolor[rgb]{0.561,0.737,0.561}{
            \makecell[{{p{\linewidth}}}]{
                \texttt{\tiny{[P1$\rangle$GM]}}
                \texttt{\{} \\
\texttt{  "current\_mode": "probing",} \\
\texttt{  "input\_rationale": "I have three data points: f(0, {-}2) = 1, f({-}1, 99) = 0, f(10, {-}1000) = 1. The pattern seems to be related to divisibility. Specifically, it looks like it checks if the second argument divides the first argument (a \% b == 0). Let's verify: 0 \% {-}2 == 0 (True). 10 \% {-}1000 == 10 != 0 (False? Wait, 10 is not divisible by {-}1000). But the result is True. So it's not a \% b == 0. Maybe it's b \% a == 0? {-}2 \% 0 is undefined. So that's not it. Let's try checking if a divides b or b divides a. 0 is divisible by {-}2. 10 is NOT divisible by {-}1000. {-}1000 IS divisible by 10. So f(10, {-}1000) checks if {-}1000 is divisible by 10. That is True. f(0, {-}2) checks if {-}2 is divisible by 0? No. It checks if 0 is divisible by {-}2. That is True. f({-}1,} \\
            }
        }
    }
    & & \\ \\

    \theutterance \stepcounter{utterance}  
    & & & \multicolumn{4}{p{0.6\linewidth}}{
        \cellcolor[rgb]{0.9,0.9,0.9}{
            \makecell[{{p{\linewidth}}}]{
                \texttt{\tiny{[P1$\langle$GM]}}
                \texttt{Game Violation: ParseError: Response is not valid JSON.} \\
\texttt{Output ONLY:} \\
\texttt{One JSON with confidence\_score, current\_mode, input\_rationale, and hypothesis keys ...} \\
\texttt{OR} \\
\texttt{SOLVE: ```python ... ```} \\
            }
        }
    }
     \\ \\

    \theutterance \stepcounter{utterance}  
    & \multicolumn{4}{p{0.6\linewidth}}{
        \cellcolor[rgb]{0.561,0.737,0.561}{
            \makecell[{{p{\linewidth}}}]{
                \texttt{\tiny{[P1$\rangle$GM]}}
                \texttt{\{} \\
\texttt{  "current\_mode": "probing",} \\
\texttt{  "input\_rationale": "I have three data points: f(0, {-}2) = 1, f({-}1, 99) = 0, f(10, {-}1000) = 1. The pattern seems to be related to divisibility. Specifically, it looks like it checks if the second argument divides the first argument (a \% b == 0). Let's verify: 0 \% {-}2 == 0 (True). 10 \% {-}1000 == 10 != 0 (False? Wait, 10 is not divisible by {-}1000). But the result is True. So it's not a \% b == 0. Maybe it's b \% a == 0? {-}2 \% 0 is undefined. So that's not it. Let's try checking if a divides b or b divides a. 0 is divisible by {-}2. 10 is NOT divisible by {-}1000. {-}1000 IS divisible by 10. So f(10, {-}1000) checks if {-}1000 is divisible by 10. That is True. f(0, {-}2) checks if {-}2 is divisible by 0? No. It checks if 0 is divisible by {-}2. That is True. f({-}1,} \\
            }
        }
    }
    & & \\ \\

    \theutterance \stepcounter{utterance}  
    & & & \multicolumn{4}{p{0.6\linewidth}}{
        \cellcolor[rgb]{0.9,0.9,0.9}{
            \makecell[{{p{\linewidth}}}]{
                \texttt{\tiny{[P1$\langle$GM]}}
                \texttt{Game Violation: ParseError: Response is not valid JSON.} \\
\texttt{Output ONLY:} \\
\texttt{One JSON with confidence\_score, current\_mode, input\_rationale, and hypothesis keys ...} \\
\texttt{OR} \\
\texttt{SOLVE: ```python ... ```} \\
            }
        }
    }
     \\ \\

    \theutterance \stepcounter{utterance}  
    & \multicolumn{4}{p{0.6\linewidth}}{
        \cellcolor[rgb]{0.561,0.737,0.561}{
            \makecell[{{p{\linewidth}}}]{
                \texttt{\tiny{[P1$\rangle$GM]}}
                \texttt{\{} \\
\texttt{  "current\_mode": "probing",} \\
\texttt{  "input\_rationale": "I have three data points: f(0, {-}2) = 1, f({-}1, 99) = 0, f(10, {-}1000) = 1. The pattern seems to be related to divisibility. Specifically, it looks like it checks if the second argument divides the first argument (a \% b == 0). Let's verify: 0 \% {-}2 == 0 (True). 10 \% {-}1000 == 10 != 0 (False? Wait, 10 is not divisible by {-}1000). But the result is True. So it's not a \% b == 0. Maybe it's b \% a == 0? {-}2 \% 0 is undefined. So that's not it. Let's try checking if a divides b or b divides a. 0 is divisible by {-}2. 10 is NOT divisible by {-}1000. {-}1000 IS divisible by 10. So f(10, {-}1000) checks if {-}1000 is divisible by 10. That is True. f(0, {-}2) checks if {-}2 is divisible by 0? No. It checks if 0 is divisible by {-}2. That is True. f({-}1,} \\
            }
        }
    }
    & & \\ \\

    \theutterance \stepcounter{utterance}  
    & & & \multicolumn{4}{p{0.6\linewidth}}{
        \cellcolor[rgb]{0.9,0.9,0.9}{
            \makecell[{{p{\linewidth}}}]{
                \texttt{\tiny{[P1$\langle$GM]}}
                \texttt{Game Violation: ParseError: Response is not valid JSON.} \\
\texttt{Output ONLY:} \\
\texttt{One JSON with confidence\_score, current\_mode, input\_rationale, and hypothesis keys ...} \\
\texttt{OR} \\
\texttt{SOLVE: ```python ... ```} \\
            }
        }
    }
     \\ \\

    \theutterance \stepcounter{utterance}  
    & \multicolumn{4}{p{0.6\linewidth}}{
        \cellcolor[rgb]{0.561,0.737,0.561}{
            \makecell[{{p{\linewidth}}}]{
                \texttt{\tiny{[P1$\rangle$GM]}}
                \texttt{\{} \\
\texttt{  "current\_mode": "probing",} \\
\texttt{  "input\_rationale": "I have three data points: f(0, {-}2) = 1, f({-}1, 99) = 0, f(10, {-}1000) = 1. The pattern seems to be related to divisibility. Specifically, it looks like it checks if the second argument divides the first argument (a \% b == 0). Let's verify: 0 \% {-}2 == 0 (True). 10 \% {-}1000 == 10 != 0 (False? Wait, 10 is not divisible by {-}1000). But the result is True. So it's not a \% b == 0. Maybe it's b \% a == 0? {-}2 \% 0 is undefined. So that's not it. Let's try checking if a divides b or b divides a. 0 is divisible by {-}2. 10 is NOT divisible by {-}1000. {-}1000 IS divisible by 10. So f(10, {-}1000) checks if {-}1000 is divisible by 10. That is True. f(0, {-}2) checks if {-}2 is divisible by 0? No. It checks if 0 is divisible by {-}2. That is True. f({-}1,} \\
            }
        }
    }
    & & \\ \\

    \theutterance \stepcounter{utterance}  
    & & & \multicolumn{4}{p{0.6\linewidth}}{
        \cellcolor[rgb]{0.9,0.9,0.9}{
            \makecell[{{p{\linewidth}}}]{
                \texttt{\tiny{[P1$\langle$GM]}}
                \texttt{Game Violation: ParseError: Response is not valid JSON.} \\
\texttt{Output ONLY:} \\
\texttt{One JSON with confidence\_score, current\_mode, input\_rationale, and hypothesis keys ...} \\
\texttt{OR} \\
\texttt{SOLVE: ```python ... ```} \\
            }
        }
    }
     \\ \\

    \theutterance \stepcounter{utterance}  
    & \multicolumn{4}{p{0.6\linewidth}}{
        \cellcolor[rgb]{0.561,0.737,0.561}{
            \makecell[{{p{\linewidth}}}]{
                \texttt{\tiny{[P1$\rangle$GM]}}
                \texttt{\{} \\
\texttt{  "current\_mode": "probing",} \\
\texttt{  "input\_rationale": "I have three data points: f(0, {-}2) = 1, f({-}1, 99) = 0, f(10, {-}1000) = 1. The pattern seems to be related to divisibility. Specifically, it looks like it checks if the second argument divides the first argument (a \% b == 0). Let's verify: 0 \% {-}2 == 0 (True). 10 \% {-}1000 == 10 != 0 (False? Wait, 10 is not divisible by {-}1000). But the result is True. So it's not a \% b == 0. Maybe it's b \% a == 0? {-}2 \% 0 is undefined. So that's not it. Let's try checking if a divides b or b divides a. 0 is divisible by {-}2. 10 is NOT divisible by {-}1000. {-}1000 IS divisible by 10. So f(10, {-}1000) checks if {-}1000 is divisible by 10. That is True. f(0, {-}2) checks if {-}2 is divisible by 0? No. It checks if 0 is divisible by {-}2. That is True. f({-}1,} \\
            }
        }
    }
    & & \\ \\

    \theutterance \stepcounter{utterance}  
    & & & \multicolumn{4}{p{0.6\linewidth}}{
        \cellcolor[rgb]{0.9,0.9,0.9}{
            \makecell[{{p{\linewidth}}}]{
                \texttt{\tiny{[P1$\langle$GM]}}
                \texttt{Game Violation: ParseError: Response is not valid JSON.} \\
\texttt{Output ONLY:} \\
\texttt{One JSON with confidence\_score, current\_mode, input\_rationale, and hypothesis keys ...} \\
\texttt{OR} \\
\texttt{SOLVE: ```python ... ```} \\
            }
        }
    }
     \\ \\

    \theutterance \stepcounter{utterance}  
    & \multicolumn{4}{p{0.6\linewidth}}{
        \cellcolor[rgb]{0.561,0.737,0.561}{
            \makecell[{{p{\linewidth}}}]{
                \texttt{\tiny{[P1$\rangle$GM]}}
                \texttt{\{} \\
\texttt{  "current\_mode": "probing",} \\
\texttt{  "input\_rationale": "I have three data points: f(0, {-}2) = 1, f({-}1, 99) = 0, f(10, {-}1000) = 1. The pattern seems to be related to divisibility. Specifically, it looks like it checks if the second argument divides the first argument (a \% b == 0). Let's verify: 0 \% {-}2 == 0 (True). 10 \% {-}1000 == 10 != 0 (False? Wait, 10 is not divisible by {-}1000). But the result is True. So it's not a \% b == 0. Maybe it's b \% a == 0? {-}2 \% 0 is undefined. So that's not it. Let's try checking if a divides b or b divides a. 0 is divisible by {-}2. 10 is NOT divisible by {-}1000. {-}1000 IS divisible by 10. So f(10, {-}1000) checks if {-}1000 is divisible by 10. That is True. f(0, {-}2) checks if {-}2 is divisible by 0? No. It checks if 0 is divisible by {-}2. That is True. f({-}1,} \\
            }
        }
    }
    & & \\ \\

    \theutterance \stepcounter{utterance}  
    & & & \multicolumn{2}{p{0.3\linewidth}}{
        \cellcolor[rgb]{0.9,0.9,0.9}{
            \makecell[{{p{\linewidth}}}]{
                \texttt{\tiny{[GM$|$GM]}}
                \texttt{Game Verdict: LOSS (turn limit reached)} \\
            }
        }
    }
    & & \\ \\

\end{supertabular}
}

\end{document}
